% GENERATED FILE. Edit canonical lessons, then run npm run generate:books.
% canonical-source-hash: fnv1a64:2e7b26486df68a4d
% canonical-lessons: HI-C428-kamana, HI-C428-phukna, HI-C428-munafa, HI-C428-bajat, HI-C428-vasulna

\chapter{Earning and Spending}
\label{ch:hi-earning-and-spending}
\hlchaptermodality{428}

\begin{chapteropening}
\textbf{By the end of this chapter:} \emph{I can talk about what I earn and what I waste, profit and budget, and money collected.}

Complete the last lesson of chapter 428: I can talk about what I earn and what I waste, profit and budget, and money collected.
\end{chapteropening}

\section[\texorpdfstring{\hi{कमाना} (kamānā)}{कमाना (kamānā)}]{\hi{कमाना} (kamānā) — to earn}

\label{lesson:HI-C428-kamana}



\begin{quote}
Before the new one: say the Hindi for a broker, then the Hindi for to pay off.
\end{quote}

\subsection*{You'll want to know: \hi{कमाना}}
\textbf{\hi{कमाना}} — \emph{kamānā} — \textquotedblleft{}to earn\textquotedblright{}.

\begin{quote}
\hi{मेरा} \hi{भाई} \hi{दुकान} \hi{से} \hi{अच्छा} \hi{पैसा} \hi{कमाता} \hi{है।}
\end{quote}

\begin{quote}
\emph{merā bhāī dukān se acchā paisā kamātā hai.}
\end{quote}

\begin{quote}
My brother earns good money from his shop.
\end{quote}

Past, with \hi{ने} (he or she …): \hi{उसने} \hi{बहुत} \hi{पैसा} \hi{कमाया}, \emph{usne bahut paisā kamāyā}.

\begin{cousinweb}
\textbf{\hi{कमाना}} (\emph{kamānā}) goes back to \emph{*karmāpayati}, \textquotedblleft{}works, earns\textquotedblright{}, a form CDIAL reconstructs (CDIAL 2897), built on \emph{karman}, \textquotedblleft{}work\textquotedblright{}. \hi{कमाई}, earnings, is its noun.
\end{cousinweb}

\begin{grammarlens}[title={What you've built}]
One more word for getting around.
\end{grammarlens}

\subsection*{Your turn}
\begin{itemize}
\raggedright
  \item \emph{Say it:} \emph{kamānā}
  \item \emph{Say it:} \emph{kamānā}, once more
  \item \emph{Say it:} \emph{merā bhāī dukān se acchā paisā kamātā hai.} — My brother earns good money from his shop.
  \item \emph{Recall:} say the Hindi for a broker, then the Hindi for to pay off, then say \emph{kamānā} again
\end{itemize}

\begin{tcolorbox}[breakable,colback=teal!4,colframe=teal!35!black,title={Before you move on}]
What does \hi{कमाना} mean? (\textquotedblleft{}To earn\textquotedblright{}.) Say it once more.
\end{tcolorbox}
\section[\texorpdfstring{\hi{फूँकना} (phū̃knā)}{फूँकना (phū̃knā)}]{\hi{फूँकना} (phū̃knā) — to blow (on); to squander}

\label{lesson:HI-C428-phukna}



\begin{quote}
Before the new one: say the Hindi for to pay off, then the Hindi for to earn.
\end{quote}

\subsection*{You'll want to know: \hi{फूँकना}}
\textbf{\hi{फूँकना}} — \emph{phū̃knā} — \textquotedblleft{}to blow (on); to squander\textquotedblright{}.

\begin{quote}
\hi{दादी} \hi{ने} \hi{फूँककर} \hi{दीया} \hi{बुझाया।}
\end{quote}

\begin{quote}
\emph{dādī ne phū̃kkar dīyā bujhāyā.}
\end{quote}

\begin{quote}
Grandmother blew out the lamp.
\end{quote}

Past, with \hi{ने} (he or she …): \hi{उसने} \hi{अपना} \hi{पैसा} \hi{फूँक} \hi{दिया}, \emph{usne apnā paisā phū̃k diyā}.

\hi{पैसा} \hi{फूँकना}, to blow your money.

\begin{cousinweb}
\textbf{\hi{फूँकना}} (\emph{phū̃knā}) goes back to \emph{*phūtka}, \textquotedblleft{}blowing\textquotedblright{}, a form CDIAL reconstructs (CDIAL 9102), which lists Hindi \emph{phū̃knā}, \textquotedblleft{}to blow\textquotedblright{}. Money that is blown on is gone like smoke.
\end{cousinweb}

\begin{grammarlens}[title={What you've built}]
One more word for getting around.
\end{grammarlens}

\subsection*{Your turn}
\begin{itemize}
\raggedright
  \item \emph{Say it:} \emph{phū̃knā}
  \item \emph{Say it:} \emph{phū̃knā}, once more
  \item \emph{Say it:} \emph{dādī ne phū̃kkar dīyā bujhāyā.} — Grandmother blew out the lamp.
  \item \emph{Recall:} say the Hindi for to pay off, then the Hindi for to earn, then say \emph{phū̃knā} again
\end{itemize}

\begin{tcolorbox}[breakable,colback=teal!4,colframe=teal!35!black,title={Before you move on}]
What does \hi{फूँकना} mean? (\textquotedblleft{}To blow (on); to squander\textquotedblright{}.) Say it once more.
\end{tcolorbox}
\section[\texorpdfstring{\hi{मुनाफ़ा} (munāfā)}{मुनाफ़ा (munāfā)}]{\hi{मुनाफ़ा} (munāfā) — a profit (in trade)}

\label{lesson:HI-C428-munafa}



\begin{quote}
Before the new one: say the Hindi for to earn, then the Hindi for to blow.
\end{quote}

\subsection*{You'll want to know: \hi{मुनाफ़ा}}
\textbf{\hi{मुनाफ़ा}} — \emph{munāfā} — \textquotedblleft{}a profit (in trade)\textquotedblright{}.

\begin{quote}
\hi{इस} \hi{साल} \hi{दुकान} \hi{में} \hi{अच्छा} \hi{मुनाफ़ा} \hi{हुआ।}
\end{quote}

\begin{quote}
\emph{is sāl dukān meṁ acchā munāfā huā.}
\end{quote}

\begin{quote}
The shop made a good profit this year.
\end{quote}

Masculine.

\begin{cousinweb}
\textbf{\hi{मुनाफ़ा}} (\emph{munāfā}) is usually traced to Arabic \emph{manāfiʿ}, \textquotedblleft{}profits, uses\textquotedblright{}, through Persian: what a trader gains over the cost.
\end{cousinweb}

\begin{grammarlens}[title={What you've built}]
One more word for getting around.
\end{grammarlens}

\subsection*{Your turn}
\begin{itemize}
\raggedright
  \item \emph{Say it:} \emph{munāfā}
  \item \emph{Say it:} \emph{munāfā}, once more
  \item \emph{Say it:} \emph{is sāl dukān meṁ acchā munāfā huā.} — The shop made a good profit this year.
  \item \emph{Recall:} say the Hindi for to earn, then the Hindi for to blow, then say \emph{munāfā} again
\end{itemize}

\begin{tcolorbox}[breakable,colback=teal!4,colframe=teal!35!black,title={Before you move on}]
What does \hi{मुनाफ़ा} mean? (\textquotedblleft{}A profit (in trade)\textquotedblright{}.) Say it once more.
\end{tcolorbox}
\section[\texorpdfstring{\hi{बजट} (bajaṭ)}{बजट (bajaṭ)}]{\hi{बजट} (bajaṭ) — a budget}

\label{lesson:HI-C428-bajat}



\begin{quote}
Before the new one: say the Hindi for to blow, then the Hindi for a profit (in trade).
\end{quote}

\subsection*{You'll want to know: \hi{बजट}}
\textbf{\hi{बजट}} — \emph{bajaṭ} — \textquotedblleft{}a budget\textquotedblright{}.

\begin{quote}
\hi{हमारा} \hi{बजट} \hi{कम} \hi{है}, \hi{इसलिए} \hi{हम} \hi{सस्ते} \hi{होटल} \hi{में} \hi{ठहरेंगे।}
\end{quote}

\begin{quote}
\emph{hamārā bajaṭ kam hai, isliye ham saste hoṭal meṁ ṭhahreṅge.}
\end{quote}

\begin{quote}
Our budget is small, so we will stay in a cheap hotel.
\end{quote}

Masculine.

\begin{cousinweb}
\textbf{\hi{बजट}} (\emph{bajaṭ}) is borrowed from English \emph{budget}, a household's or a government's.
\end{cousinweb}

\begin{grammarlens}[title={What you've built}]
One more word for getting around.
\end{grammarlens}

\subsection*{Your turn}
\begin{itemize}
\raggedright
  \item \emph{Say it:} \emph{bajaṭ}
  \item \emph{Say it:} \emph{bajaṭ}, once more
  \item \emph{Say it:} \emph{hamārā bajaṭ kam hai, isliye ham saste hoṭal meṁ ṭhahreṅge.} — Our budget is small, so we will stay in a cheap hotel.
  \item \emph{Recall:} say the Hindi for to blow, then the Hindi for a profit (in trade), then say \emph{bajaṭ} again
\end{itemize}

\begin{tcolorbox}[breakable,colback=teal!4,colframe=teal!35!black,title={Before you move on}]
What does \hi{बजट} mean? (\textquotedblleft{}A budget\textquotedblright{}.) Say it once more.
\end{tcolorbox}
\section[\texorpdfstring{\hi{वसूलना} (vasūlnā)}{वसूलना (vasūlnā)}]{\hi{वसूलना} (vasūlnā) — to collect (money owed); to recover (a debt)}

\label{lesson:HI-C428-vasulna}



\begin{quote}
Before the new one: say the Hindi for a profit (in trade), then the Hindi for a budget.
\end{quote}

\subsection*{You'll want to know: \hi{वसूलना}}
\textbf{\hi{वसूलना}} — \emph{vasūlnā} — \textquotedblleft{}to collect (money owed); to recover (a debt)\textquotedblright{}.

\begin{quote}
\hi{दुकानदार} \hi{हर} \hi{हफ़्ते} \hi{अपना} \hi{उधार} \hi{वसूलता} \hi{है।}
\end{quote}

\begin{quote}
\emph{dukāndār har hafte apnā udhār vasūltā hai.}
\end{quote}

\begin{quote}
The shopkeeper collects what he is owed every week.
\end{quote}

Past, with \hi{ने} (he or she …): \hi{उसने} \hi{अपना} \hi{क़र्ज़} \hi{वसूला}, \emph{usne apnā qarz vasūlā}.

\begin{cousinweb}
\textbf{\hi{वसूलना}} (\emph{vasūlnā}) is a Hindi verb built on \emph{vasūl}, \textquotedblleft{}received, recovered\textquotedblright{}, which is usually traced to Arabic \emph{wuṣūl}, \textquotedblleft{}arrival, receipt\textquotedblright{}, through Persian.
\end{cousinweb}

\begin{grammarlens}[title={What you've built}]
One more word for getting around.
\end{grammarlens}

\subsection*{Your turn}
\begin{itemize}
\raggedright
  \item \emph{Say it:} \emph{vasūlnā}
  \item \emph{Say it:} \emph{vasūlnā}, once more
  \item \emph{Say it:} \emph{dukāndār har hafte apnā udhār vasūltā hai.} — The shopkeeper collects what he is owed every week.
  \item \emph{Recall:} say the Hindi for a profit (in trade), then the Hindi for a budget, then say \emph{vasūlnā} again
\end{itemize}

\begin{tcolorbox}[breakable,colback=teal!4,colframe=teal!35!black,title={Before you move on}]
What does \hi{वसूलना} mean? (\textquotedblleft{}To collect (money owed); to recover (a debt)\textquotedblright{}.) Say it once more.
\end{tcolorbox}
